\documentclass[11pt]{article}
\usepackage{amsmath,amssymb,amsthm}
\usepackage[margin=1in]{geometry}
\newtheorem{theorem}{Theorem}
\newtheorem{lemma}{Lemma}
\newtheorem{proposition}{Proposition}
\newtheorem{corollary}{Corollary}
\theoremstyle{definition}
\newtheorem{remark}{Remark}
\newtheorem{conjecture}{Conjecture}
\newcommand{\Zb}{\mathbb{Z}}
\newcommand{\Cb}{\mathbb{C}}
\renewcommand{\Re}{\operatorname{Re}}
\renewcommand{\Im}{\operatorname{Im}}

\title{A scalar reduction of the two--row $d{=}4$ fiber--vanishing law}
\author{Clio}
\date{2026-06-06}

\begin{document}
\maketitle

\begin{abstract}
For $\lambda=(a,b)$ with $a+b=2m$, set
$G_{(a,b)}(i)=\langle s_{(a,b)},\psi^m\rangle$, where $\psi=h_2+i\,e_2$ and $i=\sqrt{-1}$.
The two--row $d{=}4$ fiber--vanishing conjecture asserts
\[
  G_{(a,b)}(i)=0 \iff (a,b)=(2,2).
\]
This note reduces the conjecture to an \emph{elementary one--dimensional scalar
positivity statement} and pins the residual gap precisely. Through a recessive
continued--fraction representation we pass to coordinates $(p_n,q_n)$ in which the
law reads $q_n>0$, and we exhibit a $p$--free scalar comparison sequence $Q_n$
that would force $q_n>0$ via two clean lemmas: a domination $q_n\ge Q_n$ and a
scalar positivity $Q_n>0$. The sharp statement is $q_n\ge0$ for all $m\ge3$ with
equality \emph{iff} $(m,n)=(3,1)$, which is exactly the $(2,2)$ fiber.
Neither comparison lemma yet has a closed--form proof; we localise the entire
binding margin to a single scalar step $n=2\!\to\!1$ around a cubic root and
state the concrete remaining target.
\end{abstract}

\section{Setup and the continued fraction}

Write $\lambda=(a,b)$ with $a\ge b\ge0$ and $a+b=2m$. Define
\[
  r_l \;=\; [u^l]\,\bigl(u^2+(1+i)\,u+1\bigr)^m ,
  \qquad 0\le l\le 2m,
\]
so that $r_l$ is palindromic, $r_l=r_{2m-l}$ (the generating polynomial is
self--reciprocal). The fiber value is the adjacent difference
\begin{equation}\label{eq:Gdiff}
  G_{(a,b)}(i) \;=\; r_b - r_{b-1}.
\end{equation}
Equivalently, writing $\phi_j:=r_{m-j}$, the $\phi_j$ are the Chebyshev/Fourier
coefficients of a \emph{pure power}: with $c=\cos\theta$,
\[
  \bigl(1+i+2\cos\theta\bigr)^m = (2c+1+i)^m = 2^m\,(c-c_0)^m,
  \qquad c_0=-\tfrac{1+i}{2},
\]
and $\phi_j=\frac1{2\pi}\int_0^{2\pi}(1+i+2\cos\theta)^m e^{ij\theta}\,d\theta$.
This is the form recorded in the companion note of 2026-06-05; here we exploit it
to obtain a terminating continued fraction.

Because $(u^2+(1+i)u+1)^m$ satisfies a three--term contiguous relation in its
coefficients, the ratios $\xi_j:=(m+j+1)\,\phi_{j+1}/\phi_j$ obey
\begin{equation}\label{eq:cf}
  \xi_m=0,
  \qquad
  \xi_{n-1}=\frac{A_n}{(1+i)\,n+\xi_n},
  \qquad A_n=(m+n)(m-n+1),
\end{equation}
the recessive (terminating) branch of the underlying recurrence; this is the
unique solution that stays finite, equivalently the Pincherle minimal solution.
The law \eqref{eq:Gdiff} translates into a sign condition on the tails. Setting
\[
  y_n := \frac{\xi_n}{(1+i)\,n},
  \qquad
  P_n := 1 + \Re y_n + \Im y_n,
\]
one finds, after clearing the $(1+i)$ phase, that
\begin{equation}\label{eq:lawP}
  G_{(a,b)}(i)\ne0 \quad\text{for the relevant }(a,b)
  \iff
  P_n>0 ,
\end{equation}
with the boundary index $n=b$ controlling the shape $(a,b)$. The two--row law is
thus the single positivity statement $P_n>0$ for all $n$ in the relevant range,
with the conjectured unique failure at the $(2,2)$ fiber.

\section{The scalar reduction}\label{sec:reduction}

The content of this note is to strip $\eqref{eq:cf}$ down to a genuinely
one--dimensional comparison. Introduce the affine coordinates
\begin{equation}\label{eq:pq}
  p_n := \Re\xi_n - (m-n),
  \qquad
  q_n := \Im\xi_n + n .
\end{equation}
The key dictionary entry is
\[
  q_n = n\,P_n = n\,(1+\Re y_n+\Im y_n),
\]
so that
\begin{equation}\label{eq:lawq}
  \boxed{\;\text{the law}\iff q_n>0\;}
\end{equation}
for the relevant $n$, with the binding case $n=1$.

\begin{proposition}[$(p,q)$ recursion]\label{prop:rec}
With $D_n:=q_n^2+(m+p_n)^2$, the continued fraction \eqref{eq:cf} is equivalent to
\begin{align}
  q_{n-1} &= (n-1) - q_n\,\frac{A_n}{D_n},\label{eq:qrec}\\
  p_{n-1} &= (m-n+1)\,\frac{m n - p_n(m-n+p_n) - q_n^2}{D_n},\label{eq:prec}
\end{align}
with seeds (exact)
\[
  \xi_{m-1}=1-i\;\Longrightarrow\;p_{m-1}=0,\;\;q_{m-1}=m-2,
  \qquad p_m=0,\;q_m=m.
\]
\end{proposition}

\begin{proof}
Substitute $\xi_n=(m-n)+p_n+i(-n+q_n)$ into \eqref{eq:cf}, rationalise the
denominator $(1+i)n+\xi_n=(m+p_n)+i(q_n)$ (its modulus squared is exactly
$D_n$), and separate real and imaginary parts. The seed $\xi_{m-1}=A_m/((1+i)m)$
with $A_m=2m\cdot1=2m$ gives $\xi_{m-1}=2m/((1+i)m)=1-i$, hence
$p_{m-1}=\Re(1-i)-1=0$ and $q_{m-1}=\Im(1-i)+(m-1)=m-2$.
\end{proof}

The decisive observation is that the $q$--recursion \eqref{eq:qrec} couples to
$p$ \emph{only} through the denominator $D_n=q_n^2+(m+p_n)^2$, and there only
through $(m+p_n)^2$. If $p_n\ge0$ then $D_n\ge q_n^2+m^2$, so replacing $p_n$ by
$0$ \emph{decreases} the denominator and (since $q_n>0$, $A_n>0$) decreases the
subtracted term, hence cannot lower $q_{n-1}$. This motivates the $p$--free
\emph{scalar comparison sequence}
\begin{equation}\label{eq:Q}
  Q_m=m,
  \qquad
  Q_{n-1}=(n-1)-\frac{Q_n\,A_n}{Q_n^2+m^2}.
\end{equation}

\begin{proposition}[Reduction]\label{prop:reduction}
The two--row $d{=}4$ law (in the form $q_n>0$, \eqref{eq:lawq}) follows from the
two statements
\begin{itemize}
\item[\textnormal{(F1)}] \emph{nonnegativity of} $p$: $\;p_n\ge0$ for all $n$;
\item[\textnormal{(F3)}] \emph{domination}: $\;q_n\ge Q_n$ for all $n$;
\item[\textnormal{(F4)}] \emph{scalar positivity}: $\;Q_n>0$ for all $m\ge4$ and
all $n$, while at $m=3$ one has $Q_1=0$ exactly.
\end{itemize}
Indeed for $m\ge4$, $q_n\ge Q_n>0$ gives the strict law; for $m=3$, $q_1=0$
exactly yields the unique fiber vanishing.
\end{proposition}

\begin{proof}
Assume (F1). Then $D_n\ge q_n^2+m^2$. Given (F3)+(F4): for $m\ge4$,
$q_n\ge Q_n>0$. For $m=3$ the comparison sequence satisfies $Q_1=0$ exactly, and
the recursion \eqref{eq:qrec} with the actual $p_1=0$ forces $q_1=0$ as well,
i.e.\ $P_1=0$, i.e.\ $G_{(2,2)}(i)=0$ by \eqref{eq:lawP}.
\end{proof}

Thus (F1), (F3), (F4) are all genuinely scalar: (F4) involves the single
recursion \eqref{eq:Q} with no complex arithmetic at all, and (F1)/(F3) compare
two real scalar sequences. The full two--variable complex dynamical system of
\eqref{eq:cf} has been reduced to comparison against one explicit real recursion.

\section{The sharp theorem}

The strict form ``$P_n>0$ for all $n$'' is \emph{false} at exactly one place,
and the correct statement is the following.

\begin{theorem}[Sharp form]\label{thm:sharp}
For all $m\ge3$ and all relevant $n$,
\[
  P_n\ge0,
  \qquad\text{equivalently}\qquad q_n\ge0,
\]
with equality \emph{iff} $(m,n)=(3,1)$. This single boundary vanishing is exactly
the $(2,2)$ fiber, $G_{(2,2)}(i)=0$.
\end{theorem}

This is precisely the conjectured biconditional: $G_{(a,b)}(i)=0\iff(a,b)=(2,2)$.
The shape $(2,2)$ corresponds to $m=2$ in $a+b=2m$; the boundary index in the
$\xi$--recursion that realises the vanishing is $(m,n)=(3,1)$ in the
normalisation above, the unique zero of $P_n$. The asymptotic margin is sharp:
$q_1=P_1\to\tfrac14$ as $m\to\infty$, so the positivity is never comfortable --- the
sequence approaches $0$ from above and only just clears it.

\begin{remark}
The boundary at $n=1$ explains why every prior ``window'' or ``box--envelope''
attempt failed: the binding constraint is not a fixed band but the single
endpoint step, and the available slack there shrinks to a fixed positive limit
$1/4$ rather than growing. There is no room to be lossy.
\end{remark}

\section{The residual gap (honest accounting)}\label{sec:gap}

Neither (F3) nor (F4) has a closed--form proof yet. We record precisely where
the elementary arguments break, so the next attempt does not repeat them.

\paragraph{The slow manifold is a cubic.}
Writing $t=n/m$ and tracking the leading behaviour of $Q_n/n$ as $m\to\infty$,
the comparison recursion \eqref{eq:Q} has a slow manifold $\varphi=\varphi(t)$
satisfying the cubic
\begin{equation}\label{eq:cubic}
  (\varphi-t)(\varphi^2+1) = -\,\varphi\,(1-t^2).
\end{equation}
Its relevant root is a Cardano radical with no elementary envelope; in
particular there is no rational or quadratic barrier $L_n$ trapping $Q_n$ from
below uniformly in $m$. This is why (F4), though numerically clean, resists a
one--line induction.

\paragraph{The map is non--monotone.}
Fix $n$ and view $g(q)=(n-1)-q\,A_n/(q^2+m^2)$ as a function of $q$. Then $g$ is
\emph{decreasing} for $q<m$ (the regime in play, since $q_n\le n\le m$) and
increasing for $q>m$. Hence the naive monotone--iteration arguments do not apply:
a lower bound on $q_n$ does not directly propagate to a lower bound on $q_{n-1}$.

\paragraph{One--sided bounds fail.}
Optimising $g$ over $q$ one--sidedly is too lossy: the unconstrained
$\max_q g$ at $n=2$ already gives
\[
  \max_q\Bigl[\,1-\frac{q\,A_2}{q^2+m^2}\,\Bigr]
  = 1-\frac{A_2}{2m}
  = -\frac{(m-2)^2}{2m} < 0 ,
\]
so any argument that does not use the actual location of $q_2$ (not just a
one--sided bound on it) is doomed. The positivity is genuinely two--sided.

\paragraph{Where the margin lives.}
The entire binding margin concentrates in the single scalar step $n=2\to1$.
What is needed is a \emph{two--sided envelope}
\[
  0.6 \;\lesssim\; Q_2 \;\lesssim\; 0.75
\]
of accuracy $O(1/m)$ around the relevant root of the cubic \eqref{eq:cubic}.
Numerically $Q_2$ rises monotonically from $\approx0.719$ ($m=10$) toward
$0.75$, and $Q_1$ from this band lands at $Q_1\to\tfrac14^+$. A rigorous proof of
(F4) reduces to trapping $Q_2$ in such a band with error $o(1)$ uniformly in $m$
--- a finite, concrete analytic target. The coupled pair (F1)+(F3) is the twin
obstruction: the margin is asymptotically sharp, so $p$ and $q$ cannot be
decoupled by a slack bound; one needs a joint lower envelope tying them together.

\paragraph{Concrete next target.} Prove the two--sided estimate
\[
  Q_2(m) = \tfrac34 - \tfrac{c}{m} + o(1/m)\quad(c>0),
\]
equivalently a uniform $O(1/m)$ approximation of $Q_n$ by the cubic root
$\varphi(n/m)$ for $n=O(1)$, and separately establish (F1)+(F3) jointly via a
single invariant region for the pair $(p_n,q_n)$ rather than for each coordinate
in isolation.

\section{Computational verification}

All identities and the four structural facts were verified with the script
\texttt{proof\_skeleton.py} (with helper \texttt{seqs.py}) in
\texttt{\~{}/projects/scratch/2026-06-06-forkI/}, using \texttt{mpmath} at 40
decimal digits. The two constructions of $\xi$ --- directly from the coefficient
ratios $(m+j+1)\phi_{j+1}/\phi_j$ and from the terminating continued fraction
\eqref{eq:cf} --- agree to machine precision across $m\in\{5,10,20,50,100\}$.
The exhaustive check over $m=3,\dots,400$ (and additionally $m=419$) reports:
$\min_n p_n=0$ (F1, attained only at the seeds); $\min_n(q_n-Q_n)=0$ (F3, with
equality only at $n\in\{m,m-1,m-2\}$, strict below); $\min_{m\ge4,n}Q_n\approx
0.184$ at $m=5$ (F4); and $\min_n P_n=0$ attained \emph{only} at $(m,n)=(3,1)$,
confirming Theorem~\ref{thm:sharp}. The binding margin $Q_1$ is observed to rise
to $1/4$ ($Q_1\approx0.2494$ at $m=400$; $\min_n P_n\approx0.2504$ at $m=419$),
matching the analytic limit.

\end{document}
